\documentclass{article}
\usepackage{amsmath,amssymb,xurl}
\usepackage[hidelinks]{hyperref}
% Shared commands for notes in docs/paper-gaps/.

\DeclareUnicodeCharacter{2080}{\ifmmode{}_{0}\else\textsubscript{0}\fi}
\DeclareUnicodeCharacter{2081}{\ifmmode{}_{1}\else\textsubscript{1}\fi}
\DeclareUnicodeCharacter{2082}{\ifmmode{}_{2}\else\textsubscript{2}\fi}
\DeclareUnicodeCharacter{2099}{\ifmmode{}_{n}\else\textsubscript{n}\fi}

\newcommand{\Lean}{\textsc{Lean}}
\ExplSyntaxOn
\NewDocumentCommand{\leanid}{m}{
  \ifmmode
    \textsf{\detokenize{#1}}
  \else
    \group_begin:
    \urlstyle{sf}
    \tl_set:Nn \l_tmpa_tl {#1}
    \tl_replace_all:Nnn \l_tmpa_tl {\_} {_}
    \exp_args:NV \path \l_tmpa_tl
    \group_end:
  \fi
}
\ExplSyntaxOff
\providecommand{\tr}{\operatorname{tr}}
\newcommand{\tnleanIssueBase}{https://github.com/LionSR/TNLean/issues/}
\newcommand{\tnleanPrBase}{https://github.com/LionSR/TNLean/pull/}
\newcommand{\tnleanDocBase}{https://sirui-lu.com/TNLean/docs/}
\newcommand{\ghissue}[1]{\href{\tnleanIssueBase#1}{GitHub issue \##1}}
\newcommand{\ghpr}[1]{\href{\tnleanPrBase#1}{PR \##1}}
\newcommand{\doclink}[1]{\href{\tnleanDocBase#1.html}{\path{#1}}}

% Dirac notation and the identity operator, as in blueprint/src/macros/common.tex.
% \providecommand keeps note-local definitions authoritative.
\usepackage{dsfont}
\providecommand{\ket}[1]{|#1\rangle}
\providecommand{\bra}[1]{\langle #1|}
\providecommand{\braket}[2]{\langle #1 | #2 \rangle}
\providecommand{\ketbra}[2]{|#1\rangle\!\langle #2|}
\providecommand{\Id}{\mathds{1}}
\providecommand{\one}{\mathds{1}}
\providecommand{\C}{\mathbb{C}}
\providecommand{\MN}[1]{M_{#1}(\C)}
\providecommand{\MD}{M_D(\C)}

% Verdict marker: \gapnote{<kind>}{<status>}. Kind is the policy
% classification (clarification, local-correction, scope-restriction,
% unfaithful, false-source, open-gap); status is open, wip (elimination
% underway), resolved, or historical. Machine-read by the site index;
% prints a verdict line.
\newcommand{\gapnote}[2]{\par\noindent\textbf{Verdict:} #1 (#2).\par}

\providecommand{\gapnote}[2]{\par\noindent\textbf{Verdict:} #1 (#2).\par}
\title{Physical Periodic Pair Operators in Hastings' Local Construction}
\author{}
\date{}
\begin{document}
\maketitle
\gapnote{scope-restriction}{open}

Hastings, \emph{Solving gapped Hamiltonians locally},
\href{https://arxiv.org/abs/cond-mat/0508554v3}{arXiv:cond-mat/0508554v3},
source \path{proj3new.tex}, constructs finite-range Gaussian filters at
lines 285--320 and groups them into neighbouring-supersite operators at
lines 609--639. The formalized auxiliary in
\path{TNLean/MPS/AreaLaw/LocalSolution/PhysicalPeriodicPairs.lean}
constructs this family on a regular nonempty periodic chain with $N$
blocks of exactly $2\ell$ physical sites each. It assumes a unique gapped
ground vector of the original Hamiltonian. The degenerate case of the source
and the subsequent hierarchical tensor construction are not claimed.

Let $G$ be the actual nearest-neighbour ring on the physical sites, let
$H=\sum_i h_i$, let $h_i$ act on the radius-$R$ ball about $i$, and assume
$\|h_i\|\le J$ and that every such ball has at most $S$ sites. These are the
source's finite-range hypotheses, lines 180--185. The truncated generator
for term $i$ is $\sum_{k:d_G(i,k)\le\ell-R}h_k$, as in equation (hpdef).
The approximation $M_i$ is the Gaussian filter of $h_i$ under that generator.
Its support is in the radius-$\ell$ ball, and $\|M_i\|\le J$.
The Gaussian weight includes the variance parameter as recorded in
\path{hastings05_truncated_dynamics.tex}.

The explicit error uses the actual maximum number of omitted terms whose
supports meet a truncation ball. For $m+3R\le\ell$, the distance of each such
term from the original support is at least $m$, by the triangle inequality.
No local approximation, residual estimate, assignment of filters, or
successful hierarchy is supplied as a hypothesis. Each coarse edge receives
exactly $2\ell$ filters, so its norm is at most $2\ell J$, its centred norm
at most $4\ell J$, and its residual is at most $2\ell$ times the proved
local residual. The edge joining the last and first blocks is included.
The result retains the shell count and the free Gaussian and time-window
parameters; it does not identify a printed characteristic length.

If $\|M-H\|\le\varepsilon$ and
$H-E_0I\ge\Delta(I-P_0)$, then
\[
 M-(E_0-\varepsilon)I\ge\Delta(I-P_0).
\]
Here $P_0$ is the original ground projector. This comparison follows from
the proved norm error; it is not an assumed gap of the approximating sum,
and does not assert that $P_0$ is its ground projector.

The remaining tasks are the degenerate ground-space basis construction,
including scalar terminal classes, the hierarchy applied to this physical
family, and a quantitative tensor representation of its final space.
The first is recorded in \path{hastings05_two_dim_ground_basis.tex}; the
closing-link issue in the second is recorded in
\path{hastings05_hierarchical_energy.tex}. These are separate from the
physical pair-family auxiliary.
\end{document}
